\section{Introduction}\label{sec:introduction}

Regular expressions were introduced by Stephen Kleene in the 1950s~\cite{Kleene1956}
as a notation for describing the patterns that a finite-state machine can recognise. The idea was to
capture, with a few simple operators for union, concatenation and Kleene star,
exactly the languages accepted by finite automata. That correspondence between a 
piece of syntax and a machine has proved to be one of the most durable ideas
in computer science. Regular expressions are interesting precisely because they sit
at this meeting point of theory and practice. On the theoretical side they describe a
well-understood class of languages and come with a rich algebra. On the practical
side they are
everywhere, from the lexical analysers inside every compiler, through search and
text processing, to validation of input and scanning of network traffic. A notation
that a programmer types in seconds is backed by decades of automata theory, and that
is what makes getting its foundations right worthwhile.

The common way to run a regular expression is to compile it into a finite automaton
and then let the machine consume the input, using one of the classical
constructions of Thompson~\cite{Thompson1968}, Glushkov~\cite{Glushkov1962}, or
McNaughton and Yamada~\cite{McNaughtonYamada1960}. There is also an older and more
direct idea, due to Brzozowski~\cite{B64}, that this thesis takes as its starting
point. Rather than building a separate machine, one can differentiate the expression
itself: the derivative of an expression with respect to a symbol describes what
remains to be matched after that symbol has been consumed, so iterating it across an
input word is already a way of running the expression. When derivatives are kept in
a canonical form there are only finitely many of them, and taking those canonical
derivatives as states yields a deterministic automaton with no auxiliary
construction at all. The approach has aged well and was reexamined from a modern
implementation perspective by Owens, Reppy and Turon~\cite{OwensReppyTuron}.

Exact matching, however, is often stricter than an application really wants. Text is
noisy, symbols are mistyped or substituted, and two inputs that a person would treat
as the same can differ in a handful of places. This is why fuzziness is interesting:
a great deal of real matching is approximate, and a theory that insists on equality
of symbols cannot describe it. A natural way to relax the rigidity is to replace
equality between symbols by a notion of similarity, so that two symbols match
not only when they are identical but also when they are close enough. Closeness is
measured by a value in $[0,1]$ and combined using triangular norms 
of fuzzy logic~\cite{met05}. The same idea appears in similarity-based logic
programming~\cite{Sessa2002}, where unification is replaced by matching up to a
similarity relation. This thesis asks what becomes of Brzozowski's derivative 
theory once symbols match by similarity. If the consumed symbol only has to be 
similar to the first symbol of a word, up to a threshold $\mu$, does the
derivative-based theory of regular languages and automata still hold together?

It does, for a simple reason.
When similarity is combined with the minimum triangular norm, cutting the similarity
relation at a threshold $\mu$ gives an equivalence relation. A fuzzy derivative over
the alphabet is then just an ordinary derivative over the quotient alphabet, where
symbols similar within $\mu$ are identified. The crisp theory transfers to the fuzzy
setting through this quotient construction. Regularity, the correspondence between
the syntactic and the semantic derivative, finiteness of the set of derivatives, and
correctness of the derivative automaton all carry over. The fuzzy automaton accepts
exactly the words that lie within tolerance $\mu$ of some word of the language.

Derivative and finiteness arguments are easy to get subtly wrong. A proof assistant
checks every step mechanically. The crisp theory is ergo mechanised in the Rocq proof
assistant~\cite{Rocq}, building on the Coquand--Siles formalisation of
regular-expression equivalence via Brzozowski
derivatives~\cite{CoquandRegexEquality,RocqRegexpBrzozowski}. It covers
regular-expression syntax and semantics, nullability, canonicalisation, the
derivative correspondence and automaton correctness, and the similarity
infrastructure on which the fuzzy layer rests. The full development is available in
a repository~\cite{OurRocqRepo}, with the code presented next to the
mathematics it verifies.

The rest of the thesis is organised as follows. Section~\ref{sec:literature} reviews
the existing literature on regular expressions, derivatives, their formalisation, and
fuzzy automata.
Section~\ref{sec:methodology} contains the technical development in two layers: the
crisp theory of derivatives and the derivative automaton, presented together with
its Rocq formalisation, followed by the fuzzy layer with its definitions, the
quotient construction, and the fuzzy automaton. Section~\ref{sec:conclusion}
summarises the results.
